\section{Introduction}

Earth observation is entering a representation-first stage. Open satellite constellations, long-running public archives, and large pretraining corpora have accelerated the development of remote sensing foundation models, while recent surveys have documented a rapid shift from vision-only pretraining pipelines toward multimodal, generalist, and task-adaptive Earth observation systems \citep{survey1,survey3,survey4}. By 2025, models such as SkySense++ were already reporting consistent gains across classification, detection, and segmentation on 12 Earth observation tasks in 7 domains, illustrating how quickly the field is moving from narrow benchmark models toward reusable geospatial representations \citep{skysensepp}. At the same time, new productized representations such as GSE and TESSERA suggest that the field is no longer defined only by trainable checkpoints; it is increasingly also defined by reusable embedding layers that can be queried, indexed, and redistributed as geospatial products \citep{gse,tessera}.

This recent surge did not emerge in isolation. Satellite embeddings inherit a longer remote sensing representation-learning trajectory that begins well before the current foundation-model vocabulary. Earlier Earth observation pipelines were dominated by hand-crafted spectral indices, texture descriptors, bag-of-visual-words features, and shallow classifiers designed for scene classification, land-cover mapping, or target detection. The first major break came when convolutional neural networks were transferred from natural images into remote sensing, making it possible to learn hierarchical spatial features directly from high-resolution imagery rather than engineering them manually. Early CNN-based scene classification work showed that transfer learning could outperform traditional remote sensing descriptors and established a practical template for representation reuse in Earth observation \citep{castelluccio2015}. In retrospect, those CNN feature extractors were the immediate ancestors of today's satellite embeddings: users often discarded the final classifier and reused intermediate or pooled activations as generic scene descriptors.

The second break was infrastructural as much as architectural. Free and open satellite archives, especially Landsat and Sentinel, dramatically increased the scale of accessible Earth observation data; platforms such as Google Earth Engine reduced the friction of assembling geographically distributed training corpora; and commodity GPU clusters made it feasible to train larger encoders on longer temporal stacks \citep{gee2017}. These conditions changed the field from one centered on carefully curated benchmark datasets to one capable of pretraining on global or continental imagery streams. Self-supervised learning then provided the methodological bridge. Seasonal Contrast demonstrated that remote sensing models could be pretrained directly from uncurated satellite time series without dense manual annotation, while the broader success of vision transformers shifted the community toward token-based architectures that scale more naturally with patch representations, masking, and multimodal fusion \citep{seco2021,vit2021,vitrs2021}. Once these ingredients were in place, the transition from task-specific CNN encoders to reusable Earth observation representation models became much faster.

The third break was the consolidation of those ingredients into model families that look recognizably ``foundation-like'' from a remote sensing perspective. Transformer backbones enabled patch masking and reconstruction objectives, as seen in SatMAE and ScaleMAE; contrastive learning made semantic alignment and retrieval practical, as seen in RemoteCLIP; and modality-aware tokenization opened the door to sensor-flexible encoders such as Prithvi, DOFA, Galileo, and later multimodal systems \citep{satmae,scalemae,remoteclip,prithvi,dofa,galileo}. At the same time, product-serving workflows matured: instead of exposing only a checkpoint, some systems began serving annual or static embedding layers that could be sampled directly in mapping pipelines. That shift matters because it changes the unit of use. The operational object is no longer only a model to fine-tune, but also a representation product to query, compare, and fuse with labels or side information.

This evolution is important because remote sensing is not merely another image domain. Earth observation representations must absorb multi-resolution and multi-sensor inputs, spectral bands beyond RGB, repeated observations through time, strong geographic and seasonal dependence, and large shifts in coverage, cloud conditions, and annotation density. As emphasized by recent surveys, the gap between natural-image assumptions and remote sensing practice remains substantial, especially when models are transferred across sensors, regions, and tasks \citep{survey1,survey3,survey4}. A recent NASA primer makes the same point from a deployment perspective: benchmark scores alone are insufficient for model selection in Earth observation, because practitioners must also reason about pretraining data, tokenization, architectural choices, and downstream use-case fit \citep{primer_eo}. In other words, the central question is no longer only whether a remote sensing foundation model exists, but what kind of representation it exposes, under what sensing assumptions, and with what practical trade-offs.

That framing motivates the core viewpoint of this review. In operational use, practitioners rarely begin with a full training recipe or a backbone diagram. They begin with an \emph{embedding}: a pooled scene vector for classification or retrieval, a dense token grid for segmentation, an annual embedding field for wide-area mapping, or a shared latent space aligned with text and metadata. The same backbone can induce very different reusable objects depending on pooling rules, temporal summarization, modality routing, adapter design, or product distillation. Conversely, models trained with different pretraining objectives can occupy a similar practical niche if both yield compact, transferable representations. This is why contrastive image-text systems such as RemoteCLIP \citep{remoteclip}, masked autoencoder lineages such as SatMAE and SatMAE++ \citep{satmae,satmaepp}, multimodal generalists such as Prithvi and DOFA \citep{prithvi,dofa}, and precomputed global embedding products such as GSE and TESSERA \citep{gse,tessera} should be read not as isolated threads, but as members of a common satellite embedding ecosystem.

An embedding-centered review is also useful because it makes evaluation more realistic. What matters for downstream use is not only whether a model reaches a high score on one benchmark, but whether its representation is robust, spatially appropriate, temporally interpretable, and easy to deploy in real mapping workflows. These issues are becoming more visible as the field matures. REOBench, for example, reports that current Earth observation foundation models can experience substantial degradation under realistic corruptions, with performance drops ranging from less than 1\% to more than 20\% depending on the task, architecture, and perturbation type \citep{reobench}. Such evidence reinforces the need for a review that does more than list models. It must connect model families to sensing regimes, output forms, robustness expectations, and application scenarios.

This review therefore adopts an explicitly embedding-centered scope. We focus on reusable latent representations derived from satellite and closely related Earth observation data, rather than on remote sensing foundation models in the abstract. Our aim is to bridge three levels that are often separated in the literature: the model paper, the data-and-task setting, and the benchmark protocol. Concretely, the survey makes four contributions:

\begin{enumerate}[leftmargin=1.6em,label=(\roman*)]
  \item it defines a review-specific concept of \emph{satellite embedding} and proposes a unified taxonomy spanning model family, sensing modality, temporal semantics, output granularity, and deployment form;
  \item it assembles a structured dossier for the \modelcount{} canonical models currently exposed in the \texttt{rs-embed} workspace, including publication metadata, supported sensors, spatial resolution, architecture, output dimensionality, and practical deployment characteristics;
  \item it links the model landscape to application-facing use by organizing downstream tasks, outlining a consistent small-area evaluation protocol, and clarifying where different embedding families are likely to be strong or fragile;
  \item it reports a compact evidence-backed benchmark over a shared Wuhan ROI, using frozen embeddings and lightweight probes to compare classification, fraction regression, water extraction, and dense segmentation behavior across model families.
\end{enumerate}

The remainder of this paper follows that logic. Section~2 defines the concept and taxonomy of satellite embeddings. Section~3 reorganizes the model landscape into major families and provides a detailed catalog. Sections~4 and~5 move from the review taxonomy to application requirements and comparative evaluation. Sections~6 and~7 discuss open questions and conclusions. The goal throughout is to make the review useful both as a scholarly synthesis and as a practical entry point for researchers who need to select, compare, and deploy satellite embedding models.
